% ─────────────────────────────────────────────────────────────────────────────
% Supplementary: baseline ports. Matches the header notes of fastv/sparsevlm/visionzip/ifprune_internvl.py and run_test_suite.sh.
% ─────────────────────────────────────────────────────────────────────────────
\section{Baseline implementation}
\label{sec:baselines}

We ported each baseline to InternVL3.5 from its official code and kept the selection algorithm and official
hyperparameters. All methods share the prompts, the tiling, the greedy decoding, and the measurement code. Budget
$B$ always means the number of visual tokens left in the KV cache, so a method keeps $\min(B,\text{all tokens})$.
The points below are the only decisions that the official code does not determine.

\paragraph{FastV~\cite{fastv}.}
We follow the official in-place variant: at layer 2, the attention of the last token ranks the visual tokens, the top
$B$ are kept, and the sequence is actually shortened, with the original position ids retained. The official in-place
code does not support a KV cache, so we added multi-token decoding with per-layer cache lengths. The pruning
decision is unchanged.

\paragraph{SparseVLM~\cite{sparsevlm}.}
We implement the rank-based removal of the paper at the official relative depths (layers 2, 6, and 15 of 32, scaled
to the depth of the model). The paper publishes no value of $\lambda$, so for each benchmark and budget we search the
smallest $\lambda$ for which the three stages alone reach between $B$ and $1.5B$ tokens, and use the median over 32
images. This step uses only token counts, not answers; images that need no compression are excluded. Remaining
excess tokens are clipped at the final stage. An image with at most $B$ tokens is not pruned ($\lambda=0$), as for
all other methods; the official experiments use a fixed input that always exceeds the budget, so the official code has
no such case.

\paragraph{VisionZip~\cite{visionzip}.}
Dominant and contextual tokens are computed per tile in the feature layer of the vision encoder. The contextual
share of 5/32 reproduces all pairs in the tables of the paper, and we apply it to other sizes as well. The official API
takes the numbers per crop, so we split $B$ evenly over the tiles. InternVL merges $2\times2$ patches by pixel shuffle before the
projector, so tokens are selected and merged after pixel shuffle and before the projector, averaging scores and keys in
each group as in the official Qwen2.5-VL code. The CLS token serves only as the query, since InternVL drops it.

\paragraph{IF-Prune~\cite{ifprune}.}
We load the released guide model (InternVL2.5-1B with its LoRA and importance head) without retraining. It runs once on
all tiles and the question, as in the official evaluation code, and the main model keeps the top $B$ tokens at layer
2, which corresponds to a keep ratio of $B$ divided by the number of tokens. Prefill and decoding follow our FastV
path. The official experiments use InternVL2 main models, so applying the guide to InternVL3.5 is an extrapolation of
their cross-model result. Both models produce 256 tokens per tile in the same order. Memory and TTFT include the guide model.
